\chapter{Conclusion}\label{ch:conclusion}

Hydrogen exsolution from a miscible sub-Neptune interior releases heat, and the evolution models
we checked leave it out. We added energy, nucleation and cooling-path functions to \Atm{} and
applied them to the 2026 H$_2$--MgSiO$_3$ mixing model and its source simulations. This chapter
summarizes the findings and what they imply.

\section{Summary of Findings}

\subsection{Open-Source Tools}

The published model uses oxide units, whereas \Atm{} uses MgSiO$_3$ formula units. The
conversion must act on the extensive free energy, not only on the compositions. With it, our
implementation reproduces the published phase diagram with an RMS residual of \digRms{}\,K.
Without it, the metastable interval of the baseline parcel is about nine times too short. The
new energy, nucleation and cooling-path functions are tested against independent solutions. The
full \Atm{} equilibrium solver reproduces the mixing enthalpy along one published-fit parcel path
to a relative difference of $\replayDiff{}$.

\subsection{The Heat of Exsolution}

At 1--3\,wt\% hydrogen and 2--5\,GPa, exsolution releases \hbarMOneLo{}--\hbarMOneHi{}\,MJ per
kg of hydrogen in the class-M1 refits that use the enthalpy data, and
\hbarPubLo{}--\hbarPubHi{}\,MJ\,kg$^{-1}$ in the published fit. The simulated excess enthalpies
anchor this value. They also show that the published fit puts too little excess enthalpy on the
hydrogen-poor side and too much on the hydrogen-rich side. In a convecting interior with
3\,wt\% hydrogen, the refit M1 gives a heat over the first 300\,K below the binodal of
\intRatioThree{} times the sensible heat. Just below the binodal, exsolution adds up to
\intGammaThree{} times the sensible heat capacity, so the effective heat capacity rises up to
\intCeffThree{}-fold if the hydrogen leaves as pure H$_2$. For a whole planet, the heat is
\planetRadiatedLo{}--\planetRadiatedHi{}\% of the radiated energy in a published budget. This is
a scaling estimate on a thermal history from an earlier phase diagram. A consistent
planetary value needs an evolution model built on the 2026 free energy.

\subsection{Exsolution Close to Equilibrium}

The top of a convecting interior degasses into the envelope without nucleation. Away from the
envelope, homogeneous nucleation ends metastability \dlyFigOneCPubLo{}--\dlyFigOneCPubHi{}\,K
below the binodal for the published fit, and at most \whFigOneCPubHi{}\% of the hydrogen is
withheld. Young et al.\ plotted the binodal--spinodal window for the 2025 fit
\cite{young2025}. In the 2026 fit it spans \pubIntTwo{}\,K for the baseline parcel, but it is a
mean-field upper bound that cooling parcels do not reach.

\subsection{What the Source Data Constrain}

The published fit departs from its source simulations in three places: the compositions of the
two near-critical runs, the excess volume at 5\,GPa and 4000\,K, and the excess enthalpy of the
most hydrogen-poor simulated mixture ($x = 0.24$) at the same conditions. Across four model
classes and six
data treatments, the median window ranges from \IntRangeLo{} to \IntRangeHi{}\,K. A perturbation
that leaves every bulk equilibrium quantity unchanged still moves the spinodal by
\bumpFiveAbs{} to \bumpTwoAbs{}\,K between 5 and 2\,wt\% hydrogen. At the 3\,wt\% baseline the
shift is only \bumpThree{}\,K, so there the spread between model classes dominates. The spinodal
is therefore a property of the assumed functional form. The heat is better constrained, because
it rests on the simulated enthalpies of homogeneous states.

\section{Implications}

\begin{enumerate}
\item \textbf{Energy budgets.} Evolution models with miscible interiors should include the heat
of exsolution. It should come from the same free energy as the phase boundary, so that the two
stay consistent.
\item \textbf{Timing of release.} Equilibrium exsolution at the binodal is a good approximation
for convecting interiors. The demographic tests proposed for young sub-Neptunes
\cite{rogers2026} do not need the metastable window as a systematic uncertainty.
\item \textbf{Fitting mixing models.} A mixing model used for energies must be fitted to the
excess enthalpies, not only to phase boundaries. A fit to phase compositions alone overpredicts
the heat, and the enthalpy data exclude it.
\end{enumerate}

\section{Future Outlook}

The per-run composition estimators, composition time series and interface widths of the source
simulations would sharpen both the heat and the kinetics. Homogeneous enthalpies below 4000\,K
would test the temperature dependence of the heat. A machine-learned potential trained on the
same functional could test the effect of cell size and map coexistence in large cells. It could
also measure the composition curvature through the structure factor. The next step for the
planet is an evolution model that uses the 2026 free energy for both the phase boundary and the
energy budget. It would turn the interior result into a planetary cooling history. Helium, iron and the reactive species of the melt are the next
components to add.
